
classe de @@M0@@ est dans @@M1@@, les trois coefficients @@M2@@ restent born\'es gr\^ace \`a 1.13.

Soit @@M3@@ un multiple de @@M4@@ qui est tr\`es ample. Alors dans les polyn\^omes de Hilbert @@M5@@ relativement \`a @@M6@@ des @@M7@@ dans @@M8@@, les trois coefficients @@M9@@ restent born\'es gr\^ace \`a 2.10. Par cons\'equent @@M10@@ est limit\'ee en vertu de 2.11\,; d'o\`u la premi\`ere partie de 3.8.

Quant \`a la deuxi\`eme partie de 3.8, la repr\'esentabilit\'e ``g\'en\'erique'' des foncteurs de Picard a \'et\'e \'etablie dans XII 1.2, ainsi que leur s\'eparation. Par un raisonnement similaire \`a celui de XII 2, utilisant 3.9 et XII 1.1, on se ram\`ene au cas @@M11@@ normal sur @@M12@@. Or, pour que le morphisme @@M13@@ soit quasi-affine, il suffit qu'il soit (s\'epar\'e et) quasi-fini (EGA IV 8.12.8), ce qui se v\'erifie sur les fibres g\'eom\'etriques, et r\'esulte donc du lemme suivant~:

\textbf{Lemme 3.11.} \emph{Soient @@M14@@ un corps alg\'ebriquement clos, @@M15@@ un @@M16@@-sch\'ema normal, projectif muni d'un faisceau ample @@M17@@, et @@M18@@ le sch\'ema des z\'eros d'une section @@M19@@ de @@M20@@. On suppose que @@M21@@ soit de dimension @@M22@@ (resp. @@M23@@). Alors le morphisme @@M24@@ (resp. le morphisme @@M25@@ (voir 4.1)) est de type fini et a un noyau @@M26@@ qui est un groupe fini (non-r\'eduit en g\'en\'eral).}

En effet, on peut \'evidemment supposer @@M27@@\,; donc @@M28@@ est r\'eguli\`ere, @@M29@@ \'etant int\`egre. Le morphisme envisag\'e est de type fini gr\^ace \`a la premi\`ere partie de 3.8 (resp. gr\^ace \`a 4.7 (iii)). Par suite le noyau @@M30@@ est un groupe alg\'ebrique, donc contenu dans @@M31@@, donc est le m\^eme
